% Reader question: how is a neutral judge selected, informed, and paid?
% Claim: the judge market has two different proof moments -- neutrality is
% earned before selection (visible exclusion, then a public draw), and
% accountability is an exposure interval after it (a bond that stays at risk
% through the dispute window, not a one-time label). Evidence: six candidates,
% three excluded for a principal/stake/model-family conflict with the graded
% party, a public draw from the remaining three, then submit/filter/select/
% blind grade/ledger/re-audit, with the bond exposed from select through
% re-audit. Must distinguish: the exclusion step (before) from the bond
% exposure interval (after); the slash consequence from the re-audit step
% itself. Grammar: an actor pool with a filter, feeding a step rail with an
% exposure interval underneath, rejected: a marketplace network graph (no
% order, no exposure) and a central judge circle (implies a chosen identity
% carries neutrality, not the checkable history). Five-second test: a reader
% says "excluded before the draw; exposed after it, until re-audit."
\begin{figure}[H]
\centering
\pdfigurehue{pdviolet}%
\begin{tikzpicture}[pd figure,x=1cm,y=1cm]
  \node[pd title,anchor=west] at (0,7.90) {eligibility is established before the grade};
  % ---- candidate pool and explicit exclusions --------------------------------
  \foreach \x/\lab in {1.0/$J_1$,2.3/$J_2$,3.6/$J_3$,4.9/$J_4$,6.2/$J_5$,7.5/$J_6$}
    \node[pd state,circle,minimum size=7.5mm,inner sep=0pt] (\lab) at (\x,6.60) {\lab};
  \foreach \x in {1.0,3.6,7.5} {
    \draw[pd breach rule] ({\x+.16},6.76) -- ({\x+.40},7.00);
    \draw[pd breach rule] ({\x+.16},7.00) -- ({\x+.40},6.76);
  }
  \node[pd breach label,anchor=west,text width=2.5cm] at (7.95,6.60)
    {excluded: shares a principal, stake, or model family with the graded party};
  % ---- funnel remaining candidates to a public selection point --------------
  \foreach \x in {2.3,4.9,6.2} \draw[pd hairline] (\x,6.18) -- (4.9,5.30);
  \node[pd focus datum] (draw) at (4.9,5.05) {};
  \node[pd focus label,anchor=west] at (5.15,5.05) {public draw from eligible pool};
  % ---- ceremony rail, steps named above it -----------------------------------
  \draw[pd rule,-{Stealth[length=1.8mm,width=1.1mm]}] (0.6,3.30) -- (10.6,3.30);
  \foreach \x/\lab in {1.3/submit,3.0/filter,4.9/select,6.6/{blind grade},8.3/ledger,9.9/{re-audit}} {
    \draw[pd tick] (\x,3.15) -- (\x,3.45);
    \node[pd label,anchor=south,align=center,text width=1.55cm] at (\x,3.55) {\lab};
  }
  \node[pd focus datum] at (3.0,3.30) {};
  \node[pd focus datum] at (4.9,3.30) {};
  \node[pd focus datum] at (6.6,3.30) {};
  \node[pd focus datum] at (8.3,3.30) {};
  \node[pd breach datum] at (9.9,3.30) {};
  % ---- bond exposure is an interval, not a one-time label --------------------
  \draw[pd breach rule,line width=1.45pt] (4.9,0.90) -- (9.9,0.90);
  \node[pd breach datum] at (4.9,0.90) {};
  \node[pd breach datum] at (9.9,0.90) {};
  \node[pd breach label,anchor=south east,text width=4.0cm,align=right] at (6.9,1.15)
    {judge bond remains exposed through the dispute window};
  \draw[pd breach arrow] (9.9,3.05) -- (9.9,1.15);
  \node[pd breach label,anchor=east,text width=1.9cm,align=right] at (9.65,2.45) {slash if overturned};
  % ---- the ledger: every step is one inspectable history ---------------------
  \draw[pd focus rule] (1.3,-0.20) -- (8.3,-0.20);
  \node[pd focus label,anchor=north,text width=9.0cm,align=center] at (4.8,-0.40)
    {receipt, exclusions, selection, grade, and audit remain one inspectable history};
\end{tikzpicture}
\caption{Before selection, exclusions and a public draw support neutrality. After selection, a bond and later re-audit support accountability. Both stages leave checkable evidence. [internal]}
\label{fig:judge-market}
\end{figure}
